\documentclass[11pt,a4paper]{article}

% ======================================================================
%  Operational Coherence Maintenance and the Quantum-Classical Boundary
%  Version 2 (July 2026). v1: December 2025 (ai.viXra:2512.0081v1).
%  Main changes (Appendix C):
%   1. v1's Theorem 1 ("proved here") used the unlinked-pairs infimum for
%      P_extra, which is vacuous (-infinity), and its Appendix A Step 4
%      subtracted two lower bounds (invalid). v2: the salvageable part of
%      that appendix (Steps 0-3+5) proves the unconditional entropy-
%      production bound; the coherence and extra-power bounds are imported
%      from 2512.0061 v2 (efficient baselines; free-baseline case).
%   2. The paper's own Protocol B uses a dephasing dissipator, which is
%      exactly the free-baseline case: the headline inequality is a theorem
%      with no further assumptions in the paper's own experimental setting.
%   3. RIP status updated (0064/0070/0072 v2); proxy-discipline paragraph
%      added to Protocol B (whose fixed-horizon metric and unitary-baseline
%      subtraction already comply); series references added (v1 cited none).
% ======================================================================

\usepackage[utf8]{inputenc}
\usepackage[T1]{fontenc}
\usepackage{amsmath,amssymb,amsthm,mathtools}
\usepackage[margin=1.05in]{geometry}
\usepackage[colorlinks=true,linkcolor=blue,citecolor=blue,urlcolor=blue]{hyperref}

\theoremstyle{plain}
\newtheorem{theorem}{Theorem}[section]
\newtheorem{lemma}[theorem]{Lemma}
\newtheorem{proposition}[theorem]{Proposition}
\newtheorem{corollary}[theorem]{Corollary}
\theoremstyle{definition}
\newtheorem{definition}[theorem]{Definition}
\newtheorem{assumption}[theorem]{Assumption}
\theoremstyle{remark}
\newtheorem{remark}[theorem]{Remark}

\newcommand{\Tr}{\operatorname{Tr}}
\newcommand{\kB}{k_{\mathrm{B}}}
\newcommand{\eps}{\epsilon}
\newcommand{\Ctrl}{\mathrm{Ctrl}}
\newcommand{\Pextra}{P_{\mathrm{extra}}}
\newcommand{\Pmin}{P_{\min}}
\newcommand{\Closs}{\dot{C}_{\mathrm{loss}}}

\title{Operational Coherence Maintenance and the\\
Quantum--Classical Boundary:\\
\large Formal Definitions, Falsifiable Protocols, and an Outlook for
Cognitive Systems}
\author{Lluis Eriksson\\ \small Independent Researcher\\
\small \texttt{lluiseriksson@gmail.com}}
\date{July 4, 2026 \quad (v2; v1: December 2025) \quad ai.viXra:2512.0081}

\begin{document}

\maketitle

\begin{center}
\fbox{\parbox{0.94\textwidth}{\small
\textbf{Scope note.} No collapse mechanism is proposed; the Born rule is not
derived; no claims are made about phenomenological consciousness.
\textbf{Series note (v2).} This is the integrative, protocol-facing paper of
the coherence-maintenance series (ai.viXra:2512.0060/0061/0064/0070/0071/
0072/0073, all at v2, 2026, with verification suites in a single companion
repository). v1 of this paper stated its central inequality as proved here
via a paired-strategy infimum; that formulation was vacuous, and the
corrected operational hierarchy is imported from 2512.0061 (v2)
(Theorem~\ref{thm:hierarchy}); what \emph{is} proved here
(Appendix~\ref{app:proof}) is the unconditional entropy-production bound.}}
\end{center}

\begin{abstract}
Maintaining quantum coherence against uncontrolled open-system dynamics is a
control task with unavoidable thermodynamic cost. In a finite-dimensional
setting with battery-assisted thermal operations at bath temperature $T$, we
present the corrected operational hierarchy for maintenance power: an
unconditional lower bound $\Pmin(\rho) \ge \kB T\,\sigma(\rho)$ with
$\sigma(\rho)$ the entropy production rate of the target (proved in
Appendix~\ref{app:proof}); the coherence-power bound $\Pmin \ge \kB
T\,\Closs(\rho)$ under diagonal contraction; and incremental (extra) power
bounds relative to thermodynamically efficient population-maintenance
baselines, which become assumption-free exactly when the pinched target is
stationary --- e.g.\ under pure dephasing, the very dissipator used in this
paper's numerical protocol. Here $C(\rho) = S(\rho\|\Delta[\rho])$ is
relative-entropy coherence to energy pinching and $\Closs(\rho) :=
-\frac{d}{dt}C(\rho_t)|_{t=0}$. These statements are operational,
observer-independent, and geometry-free. We then present the Rate
Inheritance Principle (RIP) as the falsifiable dynamical bridge between
static clustering and decoherence rates, with its July-2026 status: weak
form a lemma under explicit hypotheses; strong form derived, with a
frequency-resolved squared-amplitude exponent, in an exactly solvable
quasi-free local-sink class; failure through near-zero Bohr-frequency
channels realized within the secular Davies class. The exact $\omega = 0$
Dirichlet identity underlying the failure mode is quoted in its verified
form. We provide falsifiable protocols distinguishing one-shot work from
sustained maintenance power, including a numerical stress test
(distance-independent influence floor without an interface vs.\
collar-induced suppression) in a transverse-field Ising chain with remote
dephasing, whose fixed-horizon metric and unitary-baseline subtraction
already satisfy the proxy-validation discipline the series later codified.
Finally, an explicitly speculative Outlook connects the resource boundary to
the Free-Energy Principle for resource-limited agents, at a methodological
(non-phenomenological) level.
\end{abstract}

\section{Introduction}

Decoherence explains the passive suppression of interference in open quantum
systems \cite{Zurek03,Schlosshauer07,BreuerPetruccione}. This paper focuses
on a complementary, operational question: what minimal control resources are
required to \emph{prevent} coherence loss and keep a desired state stationary
against uncontrolled dynamics?

We work within an explicit thermodynamic control model (battery-assisted
thermal operations) and present the corrected hierarchy of lower bounds on
maintenance power, proved in \cite{MaintWork} (v2) and, for the unconditional
part, in Appendix~\ref{app:proof}. Separately, gapped systems exhibit static
clustering \cite{HastingsKoma,NachtergaeleSims,LiebRobinson}; whether
dynamical decoherence rates inherit that suppression is the falsifiable Rate
Inheritance Principle, whose status has moved substantially since v1
(Section~\ref{sec:rip}). Together these motivate an operational reading of
the quantum--classical boundary as a \emph{resource boundary} in state
space: coherence is ``effectively classical'' for an agent when maintaining
it exceeds the agent's power budget \cite{HCut}.

\emph{Non-claims (read first).} No collapse mechanism; no Born-rule
derivation; no claims about phenomenological consciousness. Any ``cut''
discussed here is operational (a feasibility boundary), not ontological.

\begin{remark}[What changed in v2]\label{rem:v2}
v1 defined the incremental power as an infimum over \emph{unlinked} strategy
pairs and claimed the inequality $\Pextra \ge \kB T\Closs$ as proved here.
That definition is vacuous --- the infimum is $-\infty$, since
population-maintenance strategies can burn battery at any rate --- and v1's
Appendix~A obtained the bound by subtracting two \emph{lower} bounds (its
Step~4), which is invalid; the same error pattern affected the v1 of the
companion work theorem and is documented there \cite{MaintWork}. The
corrected statements are Theorem~\ref{thm:hierarchy} below. What v1's
appendix does prove --- kept here as Appendix~\ref{app:proof} --- is the
\emph{unconditional} entropy-production bound, which needs no pairing at
all.
\end{remark}

\section{Operational framework}

\subsection{States, free energy, coherence}

$X$ finite-dimensional, $\gamma_X := e^{-\beta H_X}/\Tr e^{-\beta H_X}$,
$\beta = 1/(\kB T)$; $S(\rho\|\sigma) := \Tr[\rho(\log\rho - \log\sigma)]$
(support convention); $F_X(\rho) := \kB T\, S(\rho\|\gamma_X)$; energy
pinching $\Delta[\rho] := \sum_n\Pi_n\rho\Pi_n$; coherence $C(\rho) :=
S(\rho\|\Delta[\rho])$ (nats); loss rate $\Closs(\rho) :=
-\frac{d}{dt}C(e^{t\mathcal{L}}\rho)|_{t=0}$; entropy production
$\sigma(\rho) := -\frac{d}{dt}S(e^{t\mathcal{L}}\rho\,\|\,\gamma_S)|_{t=0}
\ge 0$ (Spohn \cite{Spohn78}), with the pinching split $\sigma =
\dot{F}^\Delta_{\mathrm{diag}} + \Closs$.

\subsection{Control model}

Strategies are stroboscopic: free evolution $e^{\Delta t\mathcal{L}}$
alternating with battery-assisted thermal operations --- global unitaries $U$
with $[U, H_S + H_B + H_W] = 0$ acting on $\rho_S \otimes \gamma_B \otimes
\sigma_W$ --- with work accounted as battery free-energy \emph{consumption}
(sign convention of \cite{MaintWork} v2). $P(s)$ is the asymptotic
consumption rate; $\Ctrl(\rho)$ the maintaining strategies.

\subsection{Incremental power, corrected definitions}

\begin{definition}[$\varepsilon$-efficient baseline \cite{MaintWork}]
$s_{\mathrm{diag}} \in \Ctrl(\Delta[\rho])$ is $\varepsilon$-efficient if
$P(s_{\mathrm{diag}}) \le \kB T\,\sigma(\Delta[\rho]) + \varepsilon$, i.e.\
it exceeds its own task's thermodynamic floor by at most $\varepsilon$.
\end{definition}

\begin{definition}[Extra power]
$\Pextra(\rho) := \Pmin(\rho) - \Pmin(\Delta[\rho])$ (difference of infima,
both over $\Ctrl$ sets assumed non-empty, the second finite).
\end{definition}

\begin{remark}[Why not unlinked pairs]\label{rem:vacuous}
$\inf_{(s, s_{\mathrm{diag}})}[P(s) - P(s_{\mathrm{diag}})] = -\infty$: fix
$s$ and let the baseline append battery-burning cycles (unitaries on $WB$
dumping battery free energy into the bath act trivially on $S$ and preserve
maintenance of $\Delta[\rho]$). v1's definition and its ``baseline
matching'' remark did not exclude this; efficiency of the baseline is the
operative constraint, and it is precisely what a careful experimental
comparison implements (Appendix~\ref{app:calibration}).
\end{remark}

\section{Main bounds: the corrected hierarchy}

\begin{theorem}[Maintenance-power hierarchy]\label{thm:hierarchy}
Under (A1)--(A4) and (A6) of v1 (finite dimensions; GKLS with thermal fixed
point; energy-conserving strokes with thermal bath inputs; battery
free-energy accounting; full-rank $\rho$):
\begin{enumerate}
\item[(a)] (\emph{Unconditional; proved in Appendix~\ref{app:proof}})
$\Pmin(\rho) \ge \kB T\,\sigma(\rho)$.
\item[(b)] (\emph{Coherence power; \cite{MaintWork} v2}) under diagonal
contraction (automatic for Davies generators), $\Pmin(\rho) \ge \kB
T\,\Closs(\rho)$.
\item[(c)] (\emph{Extra power; \cite{MaintWork} v2}) against
$\varepsilon$-efficient baselines, under population autonomy,
$P(s) - P(s_{\mathrm{diag}}) \ge \kB T\,\Closs(\rho) - \varepsilon$ for
every $s \in \Ctrl(\rho)$; and in the \emph{free-baseline} case ---
$\mathcal{L}(\Delta[\rho]) = 0$ and $\Delta[\mathcal{L}(\rho)] = 0$, e.g.\
pure dephasing --- $\Pextra(\rho) = \Pmin(\rho) \ge \kB T\,\Closs(\rho)$
with no further assumptions.
\end{enumerate}
\end{theorem}

\begin{remark}[The paper's own protocols live in the free-baseline case]
\label{rem:freebaseline}
Protocol~B below (and the numerical stress test) uses a \emph{dephasing}
dissipator. For dephasing-dominated platforms the pinched target is
stationary, the population baseline is free, and the headline inequality
$\Pextra \ge \kB T\,\Closs$ is a theorem with no baseline or contraction
assumptions --- v1's claim was thus \emph{true in its own experimental
setting} for the corrected reason.
\end{remark}

Interpretation is unchanged from v1: the bounds are operational,
observer-independent, geometry-free; $C$ is in nats, so $\kB T\,\Closs$ is a
power. v1's regularity assumption (A5) is not needed for
Theorem~\ref{thm:hierarchy} and is dropped.

\section{Static clustering and the dynamical bridge (RIP)}\label{sec:rip}

Gapped systems cluster statically \cite{HastingsKoma,NachtergaeleSims};
massive-field correlators carry Bessel envelopes $K_\nu(m\eps)$
\cite{Clustering}. Whether \emph{rates} inherit these envelopes is the Rate
Inheritance Principle, formalized in \cite{RIP} (v2): the weak
(upper-envelope) form is a lemma under explicit integrability hypotheses;
the strong (envelope-class) form is a falsifiable conjecture. Status as of
this version: \emph{derived} --- with the frequency-resolved,
squared-amplitude exponent $\kappa(\eps) \propto e^{-2q(\omega_b)\eps}$,
$\cosh q(\omega) = (\mu^2 + 4 - \omega^2)/(4\mu)$ --- in an exactly solvable
quasi-free local-sink class, verified against exact Liouvillian rapidities
\cite{HCut}; \emph{failure realized} through near-zero-frequency channels
within the secular Davies model class, whose delocalized jump operators
sustain a persistent projected rate floor \cite{Stress}.

\subsection{A concrete failure mode: near-zero-frequency floors}

For a Davies generator with coupling $S$ and an active $\omega = 0$ channel,
the exact identity
\begin{equation}
\mathcal{E}^{(0)}_\sigma(O) = \frac{\gamma(0)}{2}\,
\big\lVert [S(0), O] \big\rVert^2_{2,\sigma}
\end{equation}
(KMS norm; proved and machine-verified in \cite{DaviesCore} v2, where the
general-$\omega$ decomposition requires modular weights and is \emph{not} of
commutator form) shows that whenever $S(0)$ retains non-decaying overlap
with observables at distance $\eps$ --- e.g.\ extensive conserved components,
or the generic secular delocalization measured in \cite{Stress} --- a
distance-independent rate floor arises despite static clustering. Floors are
model-dependent (a conditional design principle, as in v1), and particularly
relevant with persistent low-frequency channels.

\section{Falsifiable protocols}\label{sec:protocols}

\textbf{Protocol A (one-shot vs sustained).} Distinguish the one-shot work
of shaping a state from the sustained power of holding it: the former scales
with free-energy differences, the latter with $\kB T\,\Closs$
(Theorem~\ref{thm:hierarchy}); geometric interfaces enter the two
differently (amplitude$^2$ vs rate; see \cite{MaintWork} \S8 and \cite{HCut}).

\textbf{Protocol B (interface geometry: floor vs collar).} As in v1, with
the local $X$-basis pinching proxy $C_x$ and the \emph{fixed-horizon}
influence metric
\begin{equation}
D_T(L) := \max_{t\in[0,T]}\big|\overline{C_x^{\mathrm{noisy}}}(t; L) -
C_x^{\mathrm{unitary}}(t; L)\big|,
\end{equation}
comparing uniform-link and collar-link geometries at matched microscopic
parameters (TFIM, $N_{\mathrm{tot}} = 14$, $h/J = 1.5$, $\gamma = 0.2$,
dephasing jump at $j_{\mathrm{diss}} = L+1$, $T = 3.0$): uniform links show
an influence floor (no suppression with $L$); collar links ($J_{\mathrm{buf}}
\ll J$) show clear suppression, with fitted rates $\mu$ increasing as the
collar opacifies (v1 Table~1 retained; parameters fully declared).

Protocol~B is an influence-suppression \emph{diagnostic}, not by itself a
direct measurement of the energy-pinching coherence appearing in
Theorem~\ref{thm:hierarchy}: its readout uses the local $X$-basis pinching
proxy $C_x$, and connecting the two requires calibrating the chosen pointer
pinching to the relevant energy/pointer sector. Its role is to test the
geometry-to-rate interface; Protocol~A tests the thermodynamic maintenance
inequality directly.

\begin{remark}[Proxy discipline (v2)]\label{rem:discipline}
The series' methodology \cite{HCut,Stress,RIP} requires rate proxies to be
validated: fixed absolute horizons (not co-moving windows), explicit
baseline subtraction, and gap-removed/geometry controls. Protocol~B already
complies with the first two by construction --- $D_T$ is fixed-horizon and
subtracts the unitary baseline --- and its uniform-vs-collar comparison
plays the role of the geometry control; a critical-point sweep ($h \to J$)
is the recommended third control for any platform implementation. This is
worth stating because the v1 evidence of two companion papers failed
precisely these checks and was withdrawn \cite{HCut,Stress}.
\end{remark}

\section{The quantum--classical boundary as a resource boundary}

Combining Theorem~\ref{thm:hierarchy} with rate behavior: where rate
inheritance holds (local-sink, sub-gap regimes), isolation is exponentially
cheap in the interface width and the cost of coherence lifetime is
logarithmic \cite{HCut}; where floors persist ($\omega \simeq 0$ channels,
within the Davies class), no separation reduces the bill and a finite budget
enforces effectively classical descriptions --- an operational cut
\cite{HCut,Closure}. Below budget, no admissible strategy can maintain,
exploit, or reverse the coherence over the relevant operational timescale.

\section{Outlook for cognitive systems (speculative)}

Unchanged in substance from v1, with the corrected anchoring. Cognitive
systems are resource-limited physical agents; the Free-Energy Principle
\cite{Friston10,ParrFriston19} concerns what representations agents should
maintain, while Theorem~\ref{thm:hierarchy} bounds what maintaining them
\emph{costs} under explicit control primitives. The conditional
multi-objective reading
\begin{equation}
\min_{q, \pi}\; \mathbb{E}_\pi\big[F_{\mathrm{var}}(q)\big] +
\lambda\, \Pextra(q, \pi), \qquad \lambda > 0,
\end{equation}
with $\Pextra$ now in its corrected sense (efficient-baseline or
free-baseline), predicts: fragile encodings that lower $F_{\mathrm{var}}$
either raise measured maintenance power or are abandoned once the budget is
capped --- a falsifiable trade-off in synthetic agents with explicit
energy/compute budgets, methodologically continuous with metabolic-cost
results in neuroscience \cite{Laughlin98,AttwellLaughlin01}. No claims about
phenomenological consciousness, nor about quantum coherence in brains.

\section{Limitations}

Theorem~\ref{thm:hierarchy} assumes finite dimensions, GKLS dynamics with
thermal fixed point, bath thermal inputs, and stroboscopic battery-assisted
thermal operations; extensions to strong coupling, non-Markovian regimes, or
algebraic (Type III) settings are conditional interfaces
\cite{MaintWork,Clustering}. RIP remains open exactly for interacting gapped
buffers with local dissipation \cite{HCut}. The cognitive Outlook is
methodological only.

\appendix

\section{Proof of Theorem~\ref{thm:hierarchy}(a)}\label{app:proof}

This is the valid core of v1's Appendix~A (its Steps 0--3 and 5; the paired
Step~4 is removed, cf.\ Remark~\ref{rem:v2}).

\emph{Step 0--1.} With $\Gamma := \gamma_S\otimes\gamma_B\otimes\gamma_W$
and $[U, H_S + H_B + H_W] = 0$, $U\Gamma U^\dagger = \Gamma$, so
$S(U\rho U^\dagger\|\Gamma) = S(\rho\|\Gamma)$ per stroke.

\emph{Step 2.} Product-reference decomposition
$S(\rho_{SBW}\|\Gamma) = \sum_X S(\rho_X\|\gamma_X) + I_\Gamma(\rho_{SBW})$
with $I_\Gamma \ge 0$; applied to input (product) and output states, this
gives the per-stroke work bound $W \ge F_S(\rho'_S) - F_S(\rho_S)$, where
$W$ is battery free-energy consumption.

\emph{Step 3.} Pinching Pythagoras: $S(\rho\|\gamma_S) =
S(\Delta[\rho]\|\gamma_S) + C(\rho)$, i.e.\ $F_S = F_S\circ\Delta + \kB T C$.

\emph{Step 5.} For any $s \in \Ctrl(\rho)$ with cycle time $\Delta t$: the
correction must implement $e^{\Delta t\mathcal{L}}(\rho) \mapsto \rho$, so by
Step~2 its per-cycle consumption obeys $W(\Delta t) \ge F_S(\rho) -
F_S(e^{\Delta t\mathcal{L}}\rho) = \kB T[S(\rho\|\gamma_S) -
S(\rho_{\Delta t}\|\gamma_S)]$. Dividing by $\Delta t$, taking $\limsup$
over cycles and then $\Delta t \to 0$:
$\Pmin(\rho) \ge \kB T\,\sigma(\rho) \ge 0$, nonnegativity by Spohn
\cite{Spohn78}. With Step~3, $\sigma = \dot{F}^\Delta_{\mathrm{diag}} +
\Closs$, from which (b) follows under diagonal contraction; (c) is proved in
\cite{MaintWork} (v2). \hfill$\square$

\section{Experimental calibration and power proxies}\label{app:calibration}

As in v1, with one clarification. Wall-plug comparisons of two strategies
``implemented with identical primitives except for the coherence-stabilizing
component'' estimate $P_{\mathrm{wall}}(s) - P_{\mathrm{wall}}
(s_{\mathrm{diag}})$; by Theorem~\ref{thm:hierarchy}(c) this lower-bounds
$\kB T\,\Closs - \varepsilon$ \emph{provided the baseline is
thermodynamically efficient} --- identical primitives is exactly the
experimental implementation of $\varepsilon$-efficiency, and in
dephasing-dominated platforms the baseline is free and the proviso is
automatic (Remark~\ref{rem:freebaseline}). Thermal anchoring and error
budgets as in v1.

\section{Changes relative to version 1}

\begin{enumerate}
\item v1's Theorem~1 and its Appendix~A proof are corrected: the
unlinked-pairs $\Pextra$ was vacuous and Step~4 subtracted lower bounds
(Remark~\ref{rem:v2}, Remark~\ref{rem:vacuous}); v2 states the corrected
hierarchy (Theorem~\ref{thm:hierarchy}), proves the unconditional part here
(Appendix~\ref{app:proof} = v1's valid Steps 0--3, 5), and imports (b), (c)
from \cite{MaintWork} (v2). v1's assumption (A5) is dropped as unnecessary.
\item Remark~\ref{rem:freebaseline}: the paper's own protocols (dephasing)
sit in the free-baseline case, where the headline inequality is
assumption-free.
\item RIP section updated to the July-2026 status (derived quasi-free with
squared exponent \cite{HCut}; Davies-class floors \cite{Stress}; lemma-level
weak form \cite{RIP}); the $\omega = 0$ identity is quoted in its verified
form with the caveat that the general-$\omega$ decomposition is \emph{not}
of commutator form \cite{DaviesCore}.
\item Proxy-discipline remark added to Protocol~B
(Remark~\ref{rem:discipline}); its fixed-horizon metric and unitary-baseline
subtraction already comply, and the critical-point control is recommended
for platform implementations. v1's Figure~1/Table~1 results are retained
with fully declared parameters.
\item Series references added (v1 cited only external literature);
calibration appendix tied to $\varepsilon$-efficient baselines; work sign
fixed to the consumption convention.
\end{enumerate}

\begin{thebibliography}{19}\small

\bibitem{Zurek03} W.~H.~Zurek, Rev. Mod. Phys. \textbf{75}, 715 (2003).
\bibitem{Schlosshauer07} M.~Schlosshauer, \emph{Decoherence and the
Quantum-to-Classical Transition} (Springer, 2007).
\bibitem{BreuerPetruccione} H.-P.~Breuer and F.~Petruccione, \emph{The
Theory of Open Quantum Systems} (OUP, 2002).
\bibitem{Davies74} E.~B.~Davies, Commun. Math. Phys. \textbf{39}, 91 (1974).
\bibitem{Spohn78} H.~Spohn, J. Math. Phys. \textbf{19}, 1227 (1978).
\bibitem{Baumgratz14} T.~Baumgratz, M.~Cramer, M.~B.~Plenio, Phys. Rev.
Lett. \textbf{113}, 140401 (2014).
\bibitem{Lostaglio15} M.~Lostaglio, D.~Jennings, T.~Rudolph, Nat. Commun.
\textbf{6}, 6383 (2015).
\bibitem{Brandao15} F.~G.~S.~L.~Brand\~ao et al., PNAS \textbf{112}, 3275
(2015).
\bibitem{HastingsKoma} M.~B.~Hastings and T.~Koma, Commun. Math. Phys.
\textbf{265}, 781 (2006).
\bibitem{NachtergaeleSims} B.~Nachtergaele and R.~Sims, Commun. Math. Phys.
\textbf{265}, 119 (2006).
\bibitem{LiebRobinson} E.~H.~Lieb and D.~W.~Robinson, Commun. Math. Phys.
\textbf{28}, 251 (1972).
\bibitem{Korzekwa16} K.~Korzekwa et al., New J. Phys. \textbf{18}, 023045
(2016).
\bibitem{Cwiklinski15} P.~\'Cwikli\'nski et al., Phys. Rev. Lett.
\textbf{115}, 210403 (2015).
\bibitem{Friston10} K.~Friston, Nat. Rev. Neurosci. \textbf{11}, 127 (2010).
\bibitem{ParrFriston19} T.~Parr and K.~Friston, Biol. Cybern. \textbf{113},
495 (2019).
\bibitem{Laughlin98} S.~B.~Laughlin et al., Nat. Neurosci. \textbf{1}, 36
(1998).
\bibitem{AttwellLaughlin01} D.~Attwell and S.~B.~Laughlin, J. Cereb. Blood
Flow Metab. \textbf{21}, 1133 (2001).

\bibitem{Clustering} L.~Eriksson, ai.viXra:2512.0060, v2 (2026; v1: 2025).
Companion repository, \texttt{papers/0060-clustering-recovery}.
\bibitem{MaintWork} L.~Eriksson, ai.viXra:2512.0061, v2 (2026; v1: 2025).
Companion repository, \texttt{papers/0061-maintenance-work}.
\bibitem{HCut} L.~Eriksson, ai.viXra:2512.0064, v2 (2026; v1: 2025).
Companion repository, \texttt{papers/0064-heisenberg-cut}.
\bibitem{Stress} L.~Eriksson, ai.viXra:2512.0070, v2 (2026; v1: 2025).
Companion repository, \texttt{papers/0070-rate-inheritance}.
\bibitem{Closure} L.~Eriksson, ai.viXra:2512.0071, v2 (2026; v1: 2025).
Companion repository, \texttt{papers/0071-program-closure}.
\bibitem{RIP} L.~Eriksson, ai.viXra:2512.0072, v2 (2026; v1: 2025).
Companion repository, \texttt{papers/0072-rip-principle}.
\bibitem{DaviesCore} L.~Eriksson, ai.viXra:2512.0073, v2 (2026; v1: 2025).
Companion repository, \texttt{papers/0073-davies-core}.

\end{thebibliography}

\end{document}
